\section{Colored balancing at a fixed shift and exceptional ordinary moments}
\label{stcol:section}

The colored Lerch framework, its open-interval uniqueness theorem, its
Mellin transform, its root-of-unity translation to generalized Stieltjes
functions, and its first exceptional moment were established in the
preserved report \emph{Gamma Determinants and Colored Stieltjes Moments}
\cite{stcol:report}.  We first supply the framework and proofs needed
here.  The additional results classify all homogeneous balances at a
\emph{fixed} shift, determine the ordinary moments associated with an
arbitrary consecutive block of vanished tails, and exhibit a golden-ratio
example at the alternating color.  They answer the dimension and moment
parts of the report's further-research question on simultaneous colored
tail zeros.  Locating all such zeros remains a separate algebraic problem.

\subsection{The colored scalar transform and its Stieltjes interpretation}

Throughout this section let
\[
 \Omega=\mathbb C\setminus[1,\infty),\qquad
 \omega\in\Omega,\qquad a>0,\qquad x>0.
\]
The branch of the Lerch function is determined by the classical integral
\cite[Eq.~25.14.5]{tw:DLMFLerch}
\begin{equation}\label{stcol:lerch}
 \Phi(\omega,s,a)
 =\frac1{\Gamma(s)}\int_0^\infty
       \frac{t^{s-1}e^{-at}}{1-\omega e^{-t}}\,dt,
 \qquad \Re s>0.
\end{equation}
For $|\omega|<1$ this agrees with
$\sum_{n\ge0}\omega^n(n+a)^{-s}$.  Subtraction of successive Taylor
terms of the density at zero, together with the zeros of $1/\Gamma(s)$,
continues it to an entire function of $s$.  Set
\begin{align}
 C_q^\omega(x,a)
  &:=\left.\partial_z^q\bigl[x^z\Phi(\omega,1+z,x+a)\bigr]\right|_{z=0}
       \label{stcol:kernel}\\
  &=\sum_{j=0}^q\binom qj(\log x)^{q-j}
                     \partial_s^j\Phi(\omega,1,x+a),\notag\\
 A_k(\omega,a)&:=\Phi(\omega,-k,a)
       =(\omega\partial_\omega+a)^k\frac1{1-\omega}.
       \label{stcol:tails}
\end{align}
There is no scalar order pole at $s=1$ in this color domain.  The extra
uncolored term $\log^{q+1}x/(q+1)$ is consequently absent from
\eqref{stcol:kernel}.

\begin{lemma}[Colored tails]\label{stcol:tail-lemma}
For every integer $K\ge q\ge0$,
\begin{equation}\label{stcol:tail}
 C_q^\omega(x,a)
 =\sum_{k=q}^K\frac{(-1)^kA_k(\omega,a)}{k!}
       \left.\partial_z^q(1+z)_k\right|_{z=0}x^{-k-1}
       +O(x^{-K-2}).
\end{equation}
The expansion is locally uniform in $(\omega,a)$, and
$C_q^\omega(x,a)=O(1+|\log x|^q)$ at zero.  For every $k\ge q$ the
Pochhammer derivative in \eqref{stcol:tail} is strictly positive.
\end{lemma}
\begin{proof}
Write $f(t)=e^{-at}/(1-\omega e^{-t})$.  Its Taylor coefficient of
$t^k$ is $(-1)^kA_k/k!$.  Taylor expansion in
\eqref{stcol:lerch}, followed by the Laplace integral, gives locally
uniformly for $z$ in a small disc
\[
 x^z\Phi(\omega,1+z,x+a)
   =\sum_{k=0}^K\frac{(-1)^kA_k}{k!}(1+z)_k x^{-k-1}
     +O(x^{-K-2}).
\]
The density denominator is bounded away from zero for positive real
$t$ on compact parameter sets, and its tail decays exponentially.
Taylor remainders therefore have the stated uniform bounds.
Cauchy differentiation on a slightly smaller disc preserves these
bounds.  Since $(1+z)_k=\prod_{h=1}^k(h+z)$ has strictly positive
coefficients in degrees $0,\ldots,k$, its derivatives have the required
positivity and vanish in higher degrees.  At zero, use the finite
sum \eqref{stcol:kernel} and analyticity at the positive shift $a$.
\end{proof}

Define $Q_q$ by the local generating identity
\begin{equation}\label{stcol:Q}
 \frac{\Gamma(u+z)}{\Gamma(u)\Gamma(1+z)}
       =\sum_{q\ge0}Q_q(u)\frac{z^q}{q!}.
\end{equation}
In particular, for every integer $m\ge1$,
\begin{equation}\label{stcol:Q-integer}
 Q_q(m)=q!\,e_q(1,1/2,\ldots,1/(m-1)),
\end{equation}
where the empty-list convention is $e_0=1$ and $e_q=0$ for $q>0$.
More generally $e_q$ vanishes when $q>m-1$.

\begin{theorem}[Colored Mellin foundation]\label{stcol:mellin}
For $q\ge0$ and $0<\Re u<q+1$,
\begin{equation}\label{stcol:mellin-formula}
 \int_0^\infty x^{u-1}C_q^\omega(x,a)\,dx
   =\frac{\pi}{\sin\pi u}\,Q_q(u)\Phi(\omega,1-u,a).
\end{equation}
The integral is ordinary and absolutely convergent.  Every apparent
pole at an integer inside this strip is removable.
\end{theorem}
\begin{proof}
Initially $0<\Re u<1$ and $z$ is sufficiently close to zero.  Absolute
Fubini applied to \eqref{stcol:lerch} gives
\[
 \int_0^\infty x^{u-1}x^z\Phi(\omega,1+z,x+a)\,dx
 =\frac{\Gamma(u+z)\Gamma(1-u)}{\Gamma(1+z)}
       \Phi(\omega,1-u,a).
\]
The zero and infinity bounds in Lemma~\ref{stcol:tail-lemma}
justify locally normal differentiation in $z$.  Equation
\eqref{stcol:Q} and Euler reflection yield
\eqref{stcol:mellin-formula} on the initial strip.  The differentiated
integral is holomorphic on $0<\Re u<q+1$ by the same bounds.
Analytic continuation proves the asserted larger strip and all its
removable integer values.
\end{proof}

For a root of unity $\omega^d=1$, $\omega\ne1$, recall the uncolored
balanced Stieltjes kernel from \eqref{stbil:kernel}, equivalently
\begin{equation}\label{stcol:uncolored}
 C_q(x,b)=\frac{\log^{q+1}x}{q+1}
       +\sum_{j=0}^q(-1)^j\binom qj
                  \gamma_j(x+b)(\log x)^{q-j}.
\end{equation}
The convention is
$\zeta(1+z,b)=z^{-1}+\sum_{j\ge0}(-1)^j\gamma_j(b)z^j/j!$.
The residue-class formula for Lerch's function gives the exact dictionary
\begin{equation}\label{stcol:dictionary}
 C_q^\omega(x,a)=\frac1d\sum_{r=0}^{d-1}
        \omega^r C_q\left(\frac xd,\frac{a+r}{d}\right).
\end{equation}
Indeed $\Phi(\omega,s,x+a)=d^{-s}\sum_r\omega^r
\zeta(s,(x+a+r)/d)$, first in its convergent series region and then
by continuation.  Because $\sum_r\omega^r=0$, both the Laurent poles
and the extra logarithmic terms in \eqref{stcol:uncolored} cancel.
Differentiating the resulting holomorphic identity proves
\eqref{stcol:dictionary} for every $q$.  Thus the ordinary integrals
below are also exact identities for finite combinations of generalized
Stieltjes functions.

\subsection{Complete homogeneous classification at one shift}

Fix $(\omega,a)$ and consider constants $d_0,\ldots,d_q$ with $d_0=1$,
allowed to depend on these fixed parameters, in
\begin{equation}\label{stcol:ansatz}
 D(x)=\sum_{\ell=0}^q d_\ell(\log x)^{q-\ell}
                  \partial_s^\ell\Phi(\omega,1,x+a).
\end{equation}
There is a unique triangular change of coordinates
\begin{align}
 D(x)&=\sum_{j=0}^q\alpha_j(\log x)^{q-j}C_j^\omega(x,a),
       \label{stcol:alpha}\\
 d_\ell&=\sum_{j=\ell}^q\binom j\ell\alpha_j,
 &\alpha_j&=\sum_{\ell=j}^q(-1)^{\ell-j}\binom\ell j d_\ell.
       \label{stcol:coordinates}
\end{align}
In these coordinates normalization is $\sum_{j=0}^q\alpha_j=1$.
These are inverse binomial transforms, obtained by expanding
\eqref{stcol:kernel}; no independence claim about numerical function
values is needed for the coordinate change.

\begin{theorem}[Fixed-shift balancing space]\label{stcol:classification}
The normalized ansatz \eqref{stcol:ansatz} satisfies the strict bound
$D(x)=O(x^{-q-1})$ as $x\to\infty$ exactly as follows.
If $A_q(\omega,a)\ne0$, its only coefficients are
$\alpha_q=1$, $\alpha_j=0$ for $j<q$.
If $A_q(\omega,a)=0$, put
\[
 t=\max\{k<q:A_k(\omega,a)\ne0\}.
\]
Then its complete coefficient set is
\begin{equation}\label{stcol:space}
 \alpha_0=\cdots=\alpha_t=0,
 \qquad \sum_{j=t+1}^q\alpha_j=1.
\end{equation}
This is an affine space of dimension $q-t-1$.  In particular, for
$q\ge1$ nonuniqueness occurs precisely when
$A_{q-1}(\omega,a)=A_q(\omega,a)=0$.  Order zero is always unique.
\end{theorem}
\begin{proof}
The integer $t$ exists when required, since $A_0=(1-\omega)^{-1}\ne0$.
Put $L=\log x$ and $P_k(z)=(1+z)_k$.  The coefficient of $x^{-k-1}$
in \eqref{stcol:alpha} is the polynomial
\begin{equation}\label{stcol:tail-polynomial}
 \frac{(-1)^k A_k}{k!}
     \sum_{j=0}^{\min(k,q)}\alpha_j P_k^{(j)}(0)L^{q-j}.
\end{equation}
For $k<q$, every nonzero $A_k$ therefore forces
$\alpha_0=\cdots=\alpha_k=0$: its distinct logarithmic powers cannot
cancel, their Pochhammer coefficients are nonzero, and a term of
order $x^{-k-1}$ cannot be absorbed by a smaller inverse power.
At $k=q$ a nonzero $A_q$ forces precisely the positive logarithmic
powers to disappear, so $\alpha_0=\cdots=\alpha_{q-1}=0$.
Normalization gives the first case.  If $A_q=0$, the last nonzero
tail below it imposes exactly the constraints through $t$.
Conversely, if these constraints hold, all surviving kernels have
$A_j=\cdots=A_q=0$ and hence $C_j^\omega=O(x^{-q-2})$.
Multiplying by the finitely many logarithmic powers in
\eqref{stcol:alpha} leaves $O(x^{-q-1})$.  This proves sufficiency,
the dimension, and the criterion for nonuniqueness.
\end{proof}

The word ``strict'' excludes $O(x^{-q-1}\log^q x)$ from the target.
Discarding only the powers below $x^{-q-1}$ would miss the boundary
logarithmic constraints and produce an incorrect larger space.
For every real $0\le\omega<1$, the series
$A_k=\sum_{n\ge0}\omega^n(n+a)^k$ is positive.  Thus all fixed-shift
orders are unique there.

There is also an algebraic description of the nonunique locus.
Writing
\begin{equation}\label{stcol:numerator}
 R_k(\omega,a)=(1-\omega)^{k+1}A_k(\omega,a),\qquad R_0=1,
\end{equation}
gives polynomials obeying
\begin{equation}\label{stcol:numerator-recurrence}
 R_{k+1}=\omega(1-\omega)\partial_\omega R_k
           +[a+(k+1-a)\omega]R_k.
\end{equation}
Since $\omega=0$ has $A_k=a^k\ne0$, the pair $A_{q-1}=A_q=0$
is equivalent in our domain to
$R_{q-1}=\partial_\omega R_{q-1}=0$.
Thus nonuniqueness corresponds to a repeated color root of this
polynomial at the fixed shift.  Theorem~\ref{stcol:classification}
determines the complete balancing space once these zeros are known;
it does not classify all positive shifts and admissible colors on
this algebraic locus.

\subsection{Consecutive vanished tails give all extra ordinary moments}

\begin{theorem}[Exceptional-moment criterion and value]
\label{stcol:exceptional}
Let $m\ge q+1$ be an integer.  The ordinary improper integral
\[
 \int_0^\infty x^{m-1}C_q^\omega(x,a)\,dx
\]
converges if and only if
\begin{equation}\label{stcol:zero-block}
 A_q(\omega,a)=A_{q+1}(\omega,a)=\cdots=A_{m-1}(\omega,a)=0.
\end{equation}
In that case it converges absolutely, and its exact value is
\begin{equation}\label{stcol:exceptional-value}
 \boxed{\displaystyle
 \int_0^\infty x^{m-1}C_q^\omega(x,a)\,dx
   =(-1)^{m+1}q!\,e_q(1,1/2,\ldots,1/(m-1))
                \partial_s\Phi(\omega,1-m,a).}
\end{equation}
The case $q=m-1$ is the single-tail exceptional identity of
\cite{stcol:report}; the criterion and coefficient above treat any
consecutive block.
\end{theorem}
\begin{proof}
At zero there is no difficulty because $m>0$.  Every Pochhammer
coefficient in \eqref{stcol:tail} is nonzero for $k\ge q$.
If the first nonzero $A_k$ between $q$ and $m-1$ occurs at $k$,
the weighted integrand has a nonzero leading constant times
$x^{m-k-2}$.  Its primitive diverges as a power, or logarithmically
when $k=m-1$; the lower-order remainder cannot cancel it.  This
proves necessity even for complex $\omega$ and ordinary conditional
convergence.  Under \eqref{stcol:zero-block} the kernel instead is
$O(x^{-m-1})$, proving absolute convergence and a holomorphic Mellin
integral on $0<\Re u<m+1$.  Continue
\eqref{stcol:mellin-formula} to a neighborhood of $u=m$ and use
\[
 \frac{\pi}{\sin\pi u}=\frac{(-1)^m}{u-m}+O(u-m),\qquad
 \Phi(\omega,1-u,a)
  =-(u-m)\partial_s\Phi(\omega,1-m,a)+O((u-m)^2).
\]
The latter expansion uses $A_{m-1}=0$.  Formula
\eqref{stcol:Q-integer} now proves
\eqref{stcol:exceptional-value}.  This is a value of an ordinary
convergent integral, rather than a prescription for a divergent one.
\end{proof}

If desired, the first nonzero $A_k$ with $k\ge q$ determines the exact
right boundary $\Re u=k+1$ of absolute Mellin convergence.  Such a
$k$ must exist: otherwise the analytic density
$e^{-at}/(1-\omega e^{-t})$ would be a polynomial near zero and hence
along the positive real axis, contradicting its nonzero value at zero
and its exponential decay at infinity.

\subsection{A golden-shift degeneracy on the unit circle}

Put $\varphi=(1+\sqrt5)/2$ and $\omega=-1$.
The Euler generating function gives
\[
 \frac{e^{-at}}{1+e^{-t}}=\frac12\sum_{k\ge0}E_k(a)\frac{(-t)^k}{k!},
 \qquad A_k(-1,a)=\frac12E_k(a).
\]
The exact factorizations
\begin{align}
 E_4(a)&=a(a-1)(a^2-a-1),\label{stcol:euler4}\\
 E_5(a)&=\tfrac12(2a-1)(a^2-a-1)^2\label{stcol:euler5}
\end{align}
show that both tails vanish at the golden shift.  The neighboring
values are
\begin{equation}\label{stcol:golden-tails}
 A_3(-1,\varphi)=\frac{\sqrt5}{8},\qquad
 A_4(-1,\varphi)=A_5(-1,\varphi)=0,\qquad
 A_6(-1,\varphi)=\frac12.
\end{equation}
Consequently the complete order-five balancing space is the line
\begin{equation}\label{stcol:golden-family}
 D_\tau(x)=(1-\tau)C_5^{-1}(x,\varphi)
           +\tau\log x\,C_4^{-1}(x,\varphi),\qquad \tau\in\mathbb C.
\end{equation}
For instance $D_{-1}=2C_5^{-1}-\log x\,C_4^{-1}$ has original
coefficients $(d_0,\ldots,d_5)=(1,6,14,16,9,2)$.
The next terms, which establish the precise end of this cancellation,
are
\begin{equation}\label{stcol:golden-leading}
 C_4^{-1}(x,\varphi)=\frac{35}{12}x^{-7}+O(x^{-8}),\qquad
 C_5^{-1}(x,\varphi)=\frac74x^{-7}+O(x^{-8}).
\end{equation}
Indeed $(1+z)_6=z^6+21z^5+175z^4+735z^3+1624z^2+1764z+720$,
and \eqref{stcol:tail} gives these constants from its fourth and fifth
derivatives divided by $1440$.

Define
$\Delta_r=\zeta'(-r,\varphi/2)-\zeta'(-r,(\varphi+1)/2)$,
where the prime differentiates the order.  The extra ordinary moments are
\begin{align}
 \int_0^\infty x^4 C_4^{-1}(x,\varphi)\,dx
     &=16\Delta_4,\label{stcol:golden45}\\
 \int_0^\infty x^5 C_4^{-1}(x,\varphi)\,dx
     &=-96\Delta_5,\label{stcol:golden46}\\
 \int_0^\infty x^5 C_5^{-1}(x,\varphi)\,dx
     &=-32\Delta_5.\label{stcol:golden56}
\end{align}
Here $Q_4(5)=1$, $Q_4(6)=3$, and $Q_5(6)=1$.
The relation
$\Phi(-1,s,a)=2^{-s}[\zeta(s,a/2)-\zeta(s,(a+1)/2)]$
converts \eqref{stcol:exceptional-value} into these formulas; its
$\log2$ term vanishes at the indicated zeros.  In particular,
\begin{equation}\label{stcol:golden-vanishing}
 \int_0^\infty x^5
       \bigl[C_4^{-1}(x,\varphi)-3C_5^{-1}(x,\varphi)\bigr]\,dx=0.
\end{equation}
Every integral here is absolutely convergent.  In explicit generalized
Stieltjes form, its kernels are
\begin{equation}\label{stcol:golden-stieltjes}
 C_q^{-1}(x,\varphi)=\frac12\sum_{j=0}^q(-1)^j\binom qj
     \left(\log\frac x2\right)^{q-j}
     \left[\gamma_j\left(\frac{x+\varphi}{2}\right)
            -\gamma_j\left(\frac{x+\varphi+1}{2}\right)\right].
\end{equation}
Thus \eqref{stcol:golden46} is a higher ordinary moment enabled by two
consecutive cancellations, while \eqref{stcol:golden45} and
\eqref{stcol:golden56} also specialize the earlier single-tail theorem.
The proof resolves the exact identities independently of numerical
quadrature; finite symbolic checks and precision-refined quadrature
provide separate reproducible diagnostics.

\subsection{The alternating classification and a finite unit-circle theorem}

The golden example has an all-order uniqueness interpretation.  The
classical theorem of Brillhart states that $E_5$ is the only Euler
polynomial with multiple zeros \cite{stcol:Brillhart}.  We use that
theorem as an external mathematical input; the finite certificates below
provide a separate, self-contained check through order eight.

\begin{corollary}[Complete alternating classification]
\label{stcol:alternating-all}
For every integer $q\ge1$ and positive real $a$,
\[
 A_{q-1}(-1,a)=A_q(-1,a)=0
 \quad\Longleftrightarrow\quad q=5,\quad a=\varphi.
\]
Consequently \eqref{stcol:golden-family} is the only nonunique
fixed-shift homogeneous balancing space at the alternating color.
Among all alternating kernels $C_q^{-1}(x,a)$, the only one with
two extra ordinary integer Mellin moments is $C_4^{-1}(x,\varphi)$;
these are its moments at $m=5,6$ in
\eqref{stcol:golden45}--\eqref{stcol:golden46}.  No such kernel has
three extra integer moments beyond the general strip $0<\Re u<q+1$.
Every other kernel has one extra moment exactly when $E_q(a)=0$,
and otherwise has none.
\end{corollary}
\begin{proof}
The generating function gives $E_q'=qE_{q-1}$, so a common zero
is a multiple zero of $E_q$.  Brillhart's theorem and
\eqref{stcol:euler5} leave only $q=5$ and the two roots of
$a^2-a-1$; precisely one of these is positive.  The nonzero
neighbors in \eqref{stcol:golden-tails} show that this zero block
has length exactly two.  The balancing and moment statements
now follow from Theorems~\ref{stcol:classification}
and~\ref{stcol:exceptional}.
\end{proof}

\begin{theorem}[Unit-circle classification through order eight]
\label{stcol:circle-eight}
For $1\le q\le8$, $a>0$, $|\omega|=1$, and $\omega\ne1$,
\begin{equation}\label{stcol:circle-locus}
 A_{q-1}(\omega,a)=A_q(\omega,a)=0
 \quad\Longleftrightarrow\quad
 (q,\omega,a)=(5,-1,\varphi).
\end{equation}
The nonalternating part of this statement is certified by sixteen
explicit rational polynomial identities.  It makes no assertion about
nonalternating colors at orders greater than eight.
\end{theorem}
\begin{proof}
Every permitted color has a unique representation
\[
 \omega=\frac{it-1}{it+1},\qquad t\in\mathbb R.
\]
Put $b=a-\tfrac12$ and $B_k=(1-\omega)A_k$.  The tail recurrence
\eqref{stcol:tails} becomes
\begin{equation}\label{stcol:circle-recurrence}
 B_0=1,\qquad
 B_{k+1}=(b+it/2)B_k+\frac{i(1+t^2)}2\partial_tB_k.
\end{equation}
For example, the coefficient of $\partial_t$ follows from
$\omega/(d\omega/dt)=i(1+t^2)/2$; differentiating the factor
$(1-\omega)^{-1}$ supplies the additional $it/2-1/2$.
With $X=b^2$, $Y=t^2$, this recurrence gives rational polynomials
$U_k,V_k$ such that
\begin{equation}\label{stcol:circle-parity}
 B_{2r}=U_{2r}(X,Y)+ibtV_{2r}(X,Y),\qquad
 B_{2r+1}=bU_{2r+1}(X,Y)+itV_{2r+1}(X,Y).
\end{equation}
Their coefficients and the certificates are supplied in
\texttt{verification/stieltjes-unit-circle-certificates.json}.

If $bt\ne0$, vanishing of both tails forces
$U_{q-1}=V_{q-1}=U_q=V_q=0$.  For every $1\le q\le8$ the
certificate gives four polynomials $h_0,h_1,h_2,h_3\in\mathbb Q[X,Y]$
whose exact identity is
\begin{equation}\label{stcol:circle-bezout}
 h_0U_{q-1}+h_1V_{q-1}+h_2U_q+h_3V_q=1.
\end{equation}
For illustration, at $q=5$ they are
\[
 (h_0,h_1,h_2,h_3)
   =\left(-\frac52,\frac{5(4X-8Y-3)}8,-\frac54,-\frac12\right).
\]
The sum in \eqref{stcol:circle-bezout} is verified by expanding
the polynomials, with no root-location approximations.
If $b=0$, $t\ne0$, the surviving component of $B_k$ is
$U_k(0,Y)$ for even $k$ and $itV_k(0,Y)$ for odd $k$.
For each consecutive pair a two-polynomial identity with sum $1$
excludes a common zero in this case as well.  The following table
records the maximum total degrees of the supplied multipliers.
\[
\begin{array}{c|rrrrrrrr}
 q&1&2&3&4&5&6&7&8\\\hline
 bt\ne0&0&0&0&0&1&1&2&2\\
 b=0,\ t\ne0&0&0&0&1&1&2&2&3
\end{array}
\]
Finally $t=0$ means $\omega=-1$.  The supplied univariate Bezout
identities and divisibilities give
\[
 \gcd(E_{q-1},E_q)=
 \begin{cases}a^2-a-1,&q=5,\\1,&q=1,2,3,4,6,7,8.\end{cases}
\]
This proves the bounded theorem independently of Brillhart's result.
The verifier reconstructs every $B_k$ directly from the original
numerators \eqref{stcol:numerator-recurrence}, then expands the
identities and checks the divisibilities.  Thus the certificate checks
also verify the change of variables used in the proof.
\end{proof}

There is also a self-contained all-order conclusion at the central shift.
Set $a=1/2$, so $b=0$ in \eqref{stcol:circle-recurrence}.  Then
\begin{equation}\label{stcol:central-polynomials}
 B_k=\left(\frac i2\right)^k P_k(t),\qquad
 P_0=1,\qquad P_{k+1}=tP_k+(1+t^2)P_k'.
\end{equation}
This identity follows immediately by induction in the same recurrence.
It also identifies $P_k$ as the classical secant derivative polynomials:
if $t=\tan v$, then
$\partial_v^k\sec v=\sec v\,P_k(t)$.
Induction shows that $P_{2j}$ has strictly positive coefficients at every
even power from zero to $2j$, and $P_{2j+1}$ at every odd power from one
to $2j+1$.  Hence no central-shift tail vanishes for real $t\ne0$,
at any order.  Equivalently, at $a=1/2$ every nonalternating unit-circle
color has unique fixed-shift balances and no extra ordinary integer
moments.  At $t=0$, precisely the odd-order tails vanish; no consecutive
pair does.  This all-order central result does not extend the bounded
noncentral assertion of Theorem~\ref{stcol:circle-eight}.

For any admissible color, the further identity
$\partial_aA_q=qA_{q-1}$ shows that the nonunique locus is also the
repeated-shift-root locus of $A_q$.  It follows first from the defining
series in $|\omega|<1$ and then from the rational tail formulas.
Together with the repeated-color-root description following
\eqref{stcol:numerator-recurrence}, this gives two algebraic approaches
to the remaining unrestricted classification.  The alternating
classification is classical; the finite unit-circle exclusion and the
Stieltjes balancing and ordinary-moment consequences are the additions
established here.
